\documentclass[a4paper]{article}

\usepackage{graphicx}
\usepackage{amsmath,amssymb}
\usepackage{minted}
\usepackage{multirow}
\usepackage{float}
\usepackage{todonotes}
\usepackage{forest}
\usepackage{xstring}
\usepackage{xcolor}
\usepackage[most]{tcolorbox}
\usepackage{xparse}
\usepackage{natbib}

\usetikzlibrary{graphs,graphdrawing}
\usegdlibrary{layered}

% for external graphviz
\input{/home/donkarlo/Dropbox/repo/nd_ai_project/src/nd_ai/knoweldge_representation/language/formal/computer/programming/tex/latex/code_clips/external_graphviz.tex}

% for tags
\input{/home/donkarlo/Dropbox/repo/nd_ai_project/src/nd_ai/knoweldge_representation/language/formal/computer/programming/tex/latex/parts/control_sequence/kind/semtag/semtag.tex}

% for links
\input{/home/donkarlo/Dropbox/repo/nd_ai_project/src/nd_ai/knoweldge_representation/language/formal/computer/programming/tex/latex/code_clips/seehere.tex}

\title{Frequency Modulated M\"obius model}
\date{}
\author{Mohammad Rahmani}

\begin{document}
    \maketitle
    
    \seeherefooter{https://journal.r-project.org/articles/RJ-2022-015/}
    
    \cite{rueda-2021-a-novel-wave-decomposition-for-oscillatory-signals}
    and
    \cite{rodriguez-collado-2021-identification-of-action-potential-waveforms}
    propose a wave decomposition method for oscillatory signals based on Frequency Modulated M\"obius (FMM) waves.
    The method represents complex oscillatory patterns as a sum of interpretable wave components with parameters controlling amplitude, phase, location, and shape.
    For action-potential-like curves, its main value is that it provides a citable mathematical basis for approximating signal waveforms by combining several FMM components.
    
    A single FMM wave can be written as follows:
    
    \begin{equation}
        X(t)
        =
        M
        +
        A
        \cos
        \Biggl[
            \beta
            +
            2
            \arctan
            \Biggl(
            \omega
            \tan
            \Bigl(
            \frac{t-\alpha}{2}
            \Bigr)
            \Biggr)
            \Biggr]
    \end{equation}
    
    \begin{itemize}
        \item $X(t)$ is the value of the FMM wave at time $t$.
        
        \item $M$ is the baseline parameter. Bound: $M \in \mathbb{R}$.
        It controls the vertical shift of the FMM signal without changing
        the waveform shape, phase, or temporal deformation.
        
        \item $t$ is time. Bound: $t \in [0, 2\pi)$.
        
        \item $A$ is the amplitude parameter. Bound: $A \geq 0$.
        It controls the vertical scale of the waveform.
        
        \item $\alpha$ is the location parameter. Bound:
        $\alpha \in [0, 2\pi)$. It controls the horizontal position of
        the waveform on the time axis.
        
        \item $\beta$ is the phase parameter. Bound:
        $\beta \in [0, 2\pi)$. It controls the internal phase of the wave.
        Together with $\alpha$ and $\omega$, it determines the temporal
        position of the maximum and minimum points of the waveform.
        
        \item $\omega$ is the shape parameter. Bound:
        $\omega \in (0, 1]$. It controls the sharpness, skewness, and
        temporal deformation of the waveform. When $\omega$ is close to
        zero, the waveform becomes more spike-like; when $\omega$ is close
        to one, the waveform becomes more similar to a sinusoidal wave.
        
        \item $\cos(\cdot)$ is the cosine function used to generate the
        oscillatory waveform.
        
        \item $\arctan(\cdot)$ is the inverse tangent function used inside
        the phase transformation.
        
        \item $\tan(\cdot)$ is the tangent function used to transform the
        time variable before applying the inverse tangent.
        
        \item $\beta + 2 \arctan\bigl(\omega
        \tan\bigl(\frac{t-\alpha}{2}\bigr)\bigr)$ is the phase function of
        the FMM wave.
    \end{itemize}
    
    
    \section{Properties}
        \begin{itemize}
            \item The period of FMM wave is $$2\pi$$
        \end{itemize}
    
    
    \section{Implementations}
        For python implementation \seeherefooter{https://github.com/FMMGroupVa/fmmpy} and for R package \seeherefooter{https://cran.r-project.org/web/packages/FMM/vignettes/FMMVignette.html}
        
        
        \bibliographystyle{plainnat}
        \bibliography{/home/donkarlo/Dropbox/repo/nd_shared_research/bibliography/bibliography.bib}
\end{document}